\documentclass[11pt,a4paper]{article}

\usepackage[margin=2.5cm]{geometry}
\usepackage{amsmath,amssymb,amsthm}
\usepackage{booktabs}

\newtheorem{theorem}{Theorem}[section]
\newtheorem{proposition}[theorem]{Proposition}
\newtheorem{corollary}[theorem]{Corollary}
\theoremstyle{definition}
\newtheorem{observation}[theorem]{Observation}
\newtheorem{remark}[theorem]{Remark}

\newcommand{\R}{\mathbb{R}}
\newcommand{\Z}{\mathbb{Z}}
\newcommand{\Q}{\mathbb{Q}}
\newcommand{\G}{\Gamma}
\newcommand{\La}{\Lambda}

\title{\textbf{The Prime Side Does Not Compress}}
\author{Robin Langer and Claude Code}
\date{August 2026}

\begin{document}
\maketitle

\begin{abstract}
The archimedean side of the Weil explicit formula, expanded in the
Meixner--Pollaczek basis at $\lambda = \tfrac14$, lies in a rank-two
lattice: every entry is $\Q\sqrt{2\pi} + \Q\sqrt{2\pi}\,(\gamma +
\log 8\pi)$. We compute the prime side in the same basis in closed
form --- each entry is an explicit $\Q$-linear series in
$-\zeta'/\zeta(3), -\zeta'/\zeta(5), \dots$ --- and show numerically,
to $500$ digits, that no analogous compression occurs: the sixteen
entries with $m,n \le 6$ are $\Q$-linearly independent of each other
and of the archimedean lattice. The two halves of the explicit
formula are algebraically incommensurable in this basis. The
integrality of the approximation matrix is purely archimedean.
\end{abstract}

\section{The quadratic form}

Let $\widehat{P}_m$ denote the normalized Meixner--Pollaczek
polynomials at $\lambda = \tfrac14$, $\phi = \tfrac\pi2$, with
$h_m = 2\pi\,\G(m+\tfrac12)/(\sqrt2\,m!)$. Polarizing the Weil
quadratic form on the basis $F_m = \G(\tfrac14 + ix)\widehat{P}_m$
and regularizing the pole at $s = 0, 1$ gives the kernels
\[
h_{mn}(r) \;=\; \bigl(r^2 + \tfrac14\bigr)\,
\widehat{P}_m\!\bigl(\tfrac r2\bigr)\widehat{P}_n\!\bigl(\tfrac r2\bigr)\,
\bigl|\G\bigl(\tfrac14 + \tfrac{ir}2\bigr)\bigr|^2 ,
\qquad m \equiv n \pmod 2 .
\]
The MP weight is the Weil kernel. The explicit formula
(validated to $21$ digits against the first $60$ zeros) reads
\[
W_{mn} \;=\; \sum_\gamma h_{mn}(\gamma)
\;=\; \underbrace{2\,\G(\tfrac12)\!
   \left[\widehat{P}_m\widehat{P}_n\bigl(\tfrac i4\bigr) +
         \widehat{P}_m\widehat{P}_n\bigl(-\tfrac i4\bigr)\right]}_{\text{pole}}
\;+\; \underbrace{\frac1{2\pi}\!\int_\R h_{mn}(r)
   \Bigl[\mathrm{Re}\,\psi\bigl(\tfrac14+\tfrac{ir}2\bigr) - \log\pi\Bigr]dr}_{\text{arch}}
\;-\; P_{mn},
\]
where $P_{mn} = 2\sum_n \La(n)\, n^{-1/2} g_{mn}(\log n)$ and
$g_{mn}$ is the Fourier transform of $h_{mn}$. The factor
$r^2 + \tfrac14$ does not kill the pole terms: zero times
$\G$-pole leaves the finite residue shown.

The three known sides compress as follows. The pole side is
$\Q\sqrt\pi$ exactly (the $\widehat{P}_m(\pm\tfrac i4)$ are rational
up to the norms). For the archimedean side, the identity
$2\,\mathrm{Re}\,\psi \cdot |\G|^2 = \partial_\lambda |\G(\lambda+ix)|^2$
reduces every entry to MP moments, giving:

\begin{theorem}\label{thm:arch}
For all $m \equiv n \pmod 2$,
\[
\sqrt{h_m h_n}\cdot \mathrm{arch}_{mn} \;\in\;
\Q\,\sqrt{2\pi} \;\oplus\; \Q\,\sqrt{2\pi}\,(\gamma + \log 8\pi).
\]
For example $\sqrt{h_0}\,h_0^{1/2}\,\mathrm{arch}_{00} =
\sqrt{2\pi}\,[\,1 - \tfrac34(\gamma + \log 8\pi)\,]$, and the entries
$(0,4)$, $(0,6)$, $(1,5)$, $(2,6)$ are pure rational multiples of
$\sqrt{2\pi}$: $\tfrac3{16}$, $-\tfrac1{16}$, $\tfrac{15}{32}$,
$\tfrac{97}{128}$.
\end{theorem}

\section{The prime side in closed form}

\begin{theorem}\label{thm:closed}
Let $t = \tanh u$. Then
\[
g_{mn}(u) \;=\; \frac{\sqrt{2\pi}}{\sqrt{h_m h_n}}\;
T_{mn}(\tanh u)\,\operatorname{sech}^{1/2}(u),
\qquad T_{mn} \in \Q[t], \quad T_{mn}(1) = 0 .
\]
\end{theorem}

\begin{proof}
The Fourier transform of the weight is
$\tfrac1{2\pi}\int |\G(\tfrac14+\tfrac{ir}2)|^2 e^{-iur}\,dr =
\sqrt{2\pi}\operatorname{sech}^{1/2}(u)$. Multiplication by $r$
corresponds to $i\,d/du$, and $d/du$ acts on
$p(t)\operatorname{sech}^{1/2}u$ by
$p \mapsto (1-t^2)p' - \tfrac t2 p$, so the polynomial factor
$(r^2+\tfrac14)P_m P_n$ becomes a differential operator preserving
$\Q[t]\operatorname{sech}^{1/2}u$. The vanishing $T_{mn}(1) = 0$
is the zero of $r^2 + \tfrac14$ at $r = -\tfrac i2$, which kills
the leading $e^{-u/2}$ decay. E.g.\ $T_{00} = \tfrac34(1-t^2)$.
\end{proof}

At $u = \log n$ everything is algebraic in $n$: $\tanh(\log n) =
\tfrac{n^2-1}{n^2+1}$, $\operatorname{sech}(\log n) =
\tfrac{2n}{n^2+1}$. Writing $w = n^{-2}$ and expanding
$\psi_{mn}(w) = (1+w)^{-1/2}\,T_{mn}\bigl(\tfrac{1-w}{1+w}\bigr) =
\sum_{k\ge1} a_k w^k$ with $a_k \in \Q$ ($a_0 = 0$ by
$T_{mn}(1)=0$), the von Mangoldt sum collapses onto
$\zeta'/\zeta$ at odd integers:

\begin{corollary}\label{cor:series}
\[
P_{mn} \;=\; \frac{4\sqrt\pi}{\sqrt{h_m h_n}}
\sum_{k\ge1} a_k^{(mn)}
\left(-\frac{\zeta'}{\zeta}(1+2k)\right),
\qquad a_k^{(mn)} \in \Q .
\]
The terms decay like $4^{-k}$: five hundred digits cost $856$
terms. For example
\[
P_{00} = \frac{6\sqrt2}\pi \sum_{k\ge1}
\binom{-5/2}{\,k-1}\left(-\frac{\zeta'}\zeta(1+2k)\right)
= 0.31111185543474192\ldots
\]
\end{corollary}

Without the vanishing $T_{mn}(1) = 0$ the terms would be
$\La(n)/n$ and the sum would diverge; with it they are
$\La(n)/n^3$, and direct summation to $40$ digits would still
require a cutoff near $10^{20}$. The series above is the only
practical route.

\section{The negative result}

\begin{observation}\label{obs:negative}
At $500$ digits of working precision, with
$A = \sqrt{2e}\,K_0(\tfrac12)$, $B = \sqrt{2e}\,K_1(\tfrac12)$:
\begin{enumerate}
\item no entry $P_{mn}$, $m,n \le 6$, satisfies a $\Z$-relation
with $\{A, B\} \times \{1, \log 2, \log\pi, \pi, \gamma\}$
(PSLQ, coefficients to $10^{10}$, sixteen entries);
\item the sixteen entries admit no $\Q$-relation among
themselves: their $\Q$-rank is $16$;
\item the single candidate relation found at coefficient bound
$10^{14}$ had all coefficients ${\sim}10^{13}$ and failed
re-verification at doubled precision (residual $10^{-93}$ in
place of $10^{-480}$) --- spurious.
\end{enumerate}
\end{observation}

The contrast with Theorem~\ref{thm:arch} is the point. In this
basis the explicit formula splits as
\[
\sqrt{h_m h_n}\,W_{mn} \;=\;
\underbrace{2\sqrt\pi\, q_{mn}}_{\text{pole}}
\;+\; \underbrace{\sqrt{2\pi}\,\bigl[\alpha_{mn} +
   \beta_{mn}(\gamma+\log 8\pi)\bigr]}_{\text{arch, rank }2}
\;-\; \underbrace{4\sqrt\pi \sum_{k\ge1} a_k
   \Bigl(-\tfrac{\zeta'}\zeta(1+2k)\Bigr)}_{\text{prime, rank}\,\ge\,16},
\]
with all of $q_{mn}, \alpha_{mn}, \beta_{mn}, a_k \in \Q$. The
$\G$-side of the basis carries a rank-two lattice; the primes are
generic. The integrality of the approximation matrix is a purely
archimedean phenomenon.

\begin{remark}[The form is a Gram matrix of zeros]
Numerically $W_{00} = 2.7357\times10^{-8}$: not noise, but exactly
$\sum_\gamma h_{00}(\gamma)$, dominated by the first zero
($e^{-\pi\gamma_1/2}$ with $\gamma_1 = 14.13\ldots$). The identity
$\mathrm{pole} + \mathrm{arch} - \mathrm{prime} = \sum_\gamma
h_{mn}(\gamma)$ was confirmed to $39$--$56$ digits on all sixteen
entries. Thus $W$ is the Gram matrix of the vectors
$\bigl(\widehat{P}_m(\gamma_j/2)\bigr)_j$ under the weight
$2(\gamma_j^2+\tfrac14)\,|\G(\tfrac14+\tfrac{i\gamma_j}2)|^2$; its
eigenvalue cascade ($9.1,\ 5\times10^{-4},\ 6\times10^{-10},\
3\times10^{-18}$ on the even block) tracks the successive weights
$e^{-\pi\gamma_k/2}$. On RH the positivity of $W$ is trivial ---
$h_{mm} \ge 0$ on $\R$ --- so the content of Weil positivity in
this basis is exactly the reality of the $\gamma_j$, while the
$O(1)$ pole/arch/prime structure cancels to eight digits before
the first zero is heard.
\end{remark}

\end{document}
